\documentclass[a4paper,11pt]{article}

\usepackage[ngerman]{babel}
\usepackage[top=3cm,bottom=3cm,left=3cm,right=3cm]{geometry} %NICHT IN KEMIE2


%\usepackage{microtype} %schöneres typesetting  mit [stretch=10] für nur 1% anstelle von 2%
%\usepackage{lmodern} %Latin Modern FONT
%\usepackage{pifont} %schönere Font



\usepackage{amsmath}
\usepackage{nicefrac}
\usepackage{amssymb}
%\usepackage{mathtools}

\usepackage{titlesec} %Um Section, etc bearbeiten zu können
\usepackage{setspace}
\usepackage{enumitem} %enumerate
\usepackage{multicol}
\usepackage[many]{tcolorbox}
\usepackage{hyperref}
\usepackage{arydshln}

%\renewcommand{\ttdefault}{lmtt} %Schriftart wird geändert zu Latin Modern Typewriter


\hypersetup{hidelinks,
			colorlinks=true,
			linkcolor=black,
			citecolor=miudiutex,
			urlcolor=miugray2
			}

\setlength{\parindent}{0pt}
\setcounter{section}{0}


\titlespacing*{\section}{0pt}{20mm}{6mm}
\titlespacing*{\subsection}{0pt}{6mm}{3mm}
\titlespacing*{\subsubsection}{0pt}{3mm}{0mm}



%\definecolor{miudiutex}{HTML}{2f3d39}
\definecolor{miudiutex}{HTML}{1d2926}
\definecolor{miugray}{HTML}{b4cccf}
\definecolor{miugray2}{HTML}{4d7365}
\definecolor{red}{HTML}{cf1f5c}
\definecolor{green}{HTML}{00A36C}
\definecolor{delphinidin}{HTML}{315794}
\definecolor{pelargonidin}{HTML}{b33807}
\definecolor{cyanidin}{HTML}{ba0746}
\definecolor{petunidin}{HTML}{b56fed}



\raggedcolumns %wichtig für Multicolumns


%================================
% SYMBOLS
%================================

\newcommand{\RA}{$\rightarrow$~}
\newcommand{\RL}{$\rightleftharpoons$~}
\newcommand{\HARP}{\ding{226}~}
%$\xrightarrow{2}$ %Pfeil mit 2 drüber


%================================
% SHORT COMMANDS (BOXES ; ETC)
%================================


\newcommand{\notiz}[1]{\begin{notizz}{}{}\begin{center}
			{\color{miugray2!50!miudiutex}#1} \end{center}\end{notizz}}
\newcommand{\zit}[2]{\begin{zitat*}{#1}{}\small #2 \end{zitat*}}

\newcommand{\frage}[2]{\begin{qst*}{#1}{}\small #2 \end{qst*}}

\newcommand{\unten}{\end{center} \tcblower \begin{center}}


\newcommand*\circled[1]{\tikz[baseline=(char.base)]{
            \node[shape=circle,draw,inner sep=0.5pt,miudiutex] (char) {\tiny #1};}}
            
\newcommand{\miusquare}{{\color{miugray2!50}\scriptsize $\blacksquare$}~}
\newcommand{\miusquarp}{{\color{petunidin!50}\scriptsize $\blacktriangleright$}~}
\newcommand{\niu}{\item[\miusquare]}
\newcommand{\nip}{\item[\miusquarp]}

%%%%% ABSAETZE

\newcommand{\pps}{\par \vspace*{1mm}}
\newcommand{\pp}{\par \vspace*{3mm}}
\newcommand{\ppbig}{\par \vspace*{8mm}}
\newcommand{\pphuge}{\par \vspace*{10mm}}

%================================
% FRAGE BOX
%================================

\tcbuselibrary{theorems,skins,hooks}
\newtcbtheorem{qst}{Frage:}
{%
	enhanced,
	breakable,
	colback = gray!15!white,
	frame hidden,
	boxrule = 0sp,
	borderline west = {2.5pt}{0pt}{red!70},
	sharp corners,
	detach title,
	before upper = \tcbtitle\par\smallskip,
	coltitle = black,
	fonttitle = \bfseries\sffamily,
	description font = \mdseries,
	separator sign none,
	segmentation style={solid, gray!50!black},
}
{th}




%================================
% ZITAT BOX
%================================

\tcbuselibrary{theorems,skins,hooks}
\newtcbtheorem{zitat}{Zitat}
{%
	enhanced,
	breakable,
	colback = gray!15!white,
	frame hidden,
	boxrule = 0sp,
	borderline west = {2.5pt}{0pt}{black!70},
	sharp corners,
	detach title,
%	before upper = \tcbtitle\par\smallskip,
	coltitle = gray!50!black,
	fonttitle = \bfseries\sffamily,
	description font = \mdseries,
%	separator sign none,
	segmentation style={solid, gray!50!black},
}
{th}



%================================
% SIMPLE SCHWARZER RAHMEN BOX
%================================


\newcommand{\boxfr}[1]{
	\begin{tcolorbox}[
	drop shadow southeast,
	enhanced,
	colframe=black!50,
	colback=white,
	boxrule = 1pt, 
%	toprule=0pt, bottomrule=0pt,rightrule=0pt,leftrule=0pt,
	boxsep=5pt,
	arc=1pt,
	]
	\begin{center}
	#1
	\end{center}
	\end{tcolorbox}
}

\definecolor{miudiutex}{HTML}{1d2926}
\definecolor{miugray}{HTML}{b4cccf}
\definecolor{miugray2}{HTML}{4d7365}
\definecolor{white2}{HTML}{f5f7f8}
\definecolor{red}{HTML}{cf1f5c}
\definecolor{green}{HTML}{00A36C}
\definecolor{delphinidin}{HTML}{315794}
\definecolor{uniblau}{HTML}{004e9f}
\definecolor{unigelb}{HTML}{fcba00}
\definecolor{unigrau}{HTML}{909085}
\definecolor{pelargonidin}{HTML}{b33807}
\definecolor{cyanidin}{HTML}{ba0746}
\definecolor{paeonidin}{HTML}{d15ca6}
\definecolor{petunidin}{HTML}{b56fed}
\definecolor{malvidin}{HTML}{8a2d8c}

\newcommand{\metermilli}{\mskip3mu \text{mm}}
\newcommand{\cm}{\mskip3mu \text{cm}}
\newcommand{\m}{ \text{m}}
\newcommand{\lite}{ \text{l}}
\newcommand{\kg}{ \text{kg}}
\newcommand{\g}{ \text{g}}
\newcommand{\km}{ \text{km}}
\newcommand{\h}{ \text{h}}
\newcommand{\s}{ \text{s}}
\newcommand{\tonne}{ \text{t}}
\usepackage[utf8]{inputenc}
\usepackage[T1]{fontenc}
\usepackage{lmodern}

\usepackage[
    activate={true,nocompatibility},
    final,
    tracking=true,
    kerning=true,
    spacing=true,
    factor=1100,
    stretch=30,
    shrink=20
]{microtype}
\usepackage{pdfpages}
\usepackage{awesomebox}
\usepackage{marginnote}
\tcbuselibrary{documentation}
\usepackage{sidenotes}
\usepackage{import}
\usepackage{xifthen}
\usepackage{transparent}
\usepackage[table]{xcolor}


\newcommand{\abbicon}{%
  \protect\tikz[baseline=0.2mm,scale=0.8]{%
    % Das blaue Quadrat mit abgerundeten Ecken
    \fill[uniblau, rounded corners=1.2pt] (0,0) rectangle (0.8em, 0.8em);
    % Der weiße Kreis in der Mitte
    \fill[white] (0.4em, 0.4em) circle (0.12em);
  }%
} 
\usepackage[labelfont={bf},font={color=miudiutex,small},figurename={\abbicon $~$Abbildung}]{caption}


\newcommand{\bckgrnd}{
\AddToHookNext{shipout/background}{
\begin{tikzpicture}[remember picture, overlay]
 \draw[draw=none,fill=uniblau!70, shift={(current page.north west)}] (-3,0) rectangle (23,-3.75);
 \draw[draw=none,fill=miugray, shift={(current page.north east)}] (-7,0) rectangle (3,-3.75);
\end{tikzpicture}
}
}



\newcommand{\incfig}[2]{%
\def\svgwidth{#2\columnwidth}
\import{C://LaTeX-Files/pngs/}{#1.pdf_tex}
}

\renewcommand{\textbf}[1]{\textsf{{\bfseries {#1}}}}



\titleformat{\section}[hang]
		{\LARGE \color{uniblau!75!black} \sffamily \bfseries}
		{\thesection}
		{6mm}
		{}
		[]
\titleformat{\subsection}[hang]
		{\large \sffamily \bfseries \color{black}}
		{\thesubsection}
		{4mm}
		{{\color{miugray}$\blacksquare ~$}}

\titleformat{\subsubsection}[block]
		{\normalsize \bfseries \sffamily \color{miudiutex}}%
		{\normalsize \color{black} \thesubsubsection}
		{2mm}
		{{\color{miugray}$\blacksquare ~$}}
		[ \color{black}]
		
\titleformat{\paragraph}
		{\normalsize \bfseries \sffamily \color{black}}
		{\footnotesize \theparagraph}
		{1mm}
		{}

\pagestyle{empty}
\titlespacing*{\section}{0pt}{20mm}{12mm}
\titlespacing{\section}{0pt}{20mm}{12mm}

\titlespacing*{\subsection}{0pt}{14mm}{4mm}
\titlespacing{\subsection}{0pt}{14mm}{4mm}

\titlespacing*{\subsubsection}{0pt}{6mm}{3mm}
\titlespacing{\subsubsection}{0pt}{6mm}{3mm}

\titlespacing*{\paragraph}{0pt}{6mm}{2mm}
\titlespacing{\paragraph}{0pt}{6mm}{2mm}


\let\oldmarginfigure\marginfigure
\let\endoldmarginfigure\endmarginfigure

\let\oldmarginfigure\marginfigure
\let\endoldmarginfigure\endmarginfigure

\renewenvironment{marginfigure}[1][0pt] % [1] steht für 1 Argument, [0pt] ist der Standardwert
  {\oldmarginfigure[#1]\captionsetup{labelfont={sf,bf,color=uniblau!75!black},font={color=miudiutex,small}}}
  {\endoldmarginfigure}
  
  
  
  
\renewcommand{\labelitemi}{\color{uniblau!70}$\bullet$}
\renewcommand{\labelitemii}{\color{uniblau!70}$\cdot$}
\newcommand{\plmtsquare}{{\color{uniblau!70}$\blacksquare~$}}
\newcommand{\splmt}[1]{{\color{uniblau!75!black} \sffamily #1}}
\newcommand{\fplmt}[1]{{\color{uniblau!75!black} \sffamily \bfseries #1}}

\newcommand{\antwort}[1]{
\begin{tcolorbox}[
	enhanced,
	enforce breakable,
	standard jigsaw,
	boxrule=0.7pt,
	colframe=black,
	sharp corners,
	boxsep=5pt,
	title=\bfseries \sffamily Antwort,
	opacityback=0,
	coltitle=miudiutex,
	attach boxed title to top text left={yshift=-\tcboxedtitleheight/2},
	boxed title style={colback=white,colframe=white},
	bottomrule at break=0pt,
	toprule at break=0pt,
	pad at break=16mm,
]
	\begin{center}
	#1
	\end{center}
	\end{tcolorbox}
}




\rowcolors{2}{miugray!50}{miugray!50}
\arrayrulecolor{white} % Weiße Gitterlinien
\setlength{\arrayrulewidth}{1.5pt} % Etwas dickere Linien
\renewcommand{\arraystretch}{1.4} % Mehr Zeilenhöhe für bessere Lesbarkeit}


%\setlist[itemize]{itemsep=0mm}
\setlist[enumerate]{itemsep=0mm}
\usepackage[table]{xcolor}
\usepackage{pgfplots}
    \pgfplotsset{
        every axis post/.style={
	yticklabel=\empty,
	xticklabel=\empty,
	y label style={at={(axis description cs:-0.2,1.1)}},
	x label style={at={(axis description cs:1.05,0.2)}},
	ticks=none,
	width=6cm,
	axis lines=middle,
%	xlabel near ticks,
        },
    }
\usetikzlibrary{positioning, arrows.meta, calc}



\newif\ifshowme
%\showmetrue




\begin{document}
\newgeometry{top=1.8cm,bottom=4cm,left=1.5cm,right=7.5cm,marginparsep=-2.00cm,marginparwidth=4.75cm}


\bckgrnd

\begin{center}
\LARGE
\color{white}
\textbf{Übung 02} \par \vspace*{2mm}
\large \color{miugray} \textbf{Prozessbezogene Lebensmitteltechnologie}
\end{center}
\par 
\vspace*{6mm}


\color{miudiutex}
\section*{Strömungslehre II}

\subsection*{Aufgabe 1}
Öl mit einer Viskosität von 0,031 kg/m$\cdot$s und einer Dichte von 850 kg/m$^3$ soll mit 13 m$^3$/h bei einem zulässigen Druckabfall von 0,1 bar laminar durch ein kreisrundes Rohr der Länge $l$ = 100 m geleitet werden. \par 
\begin{itemize}
\item Wie groß ist die Nennweite?
\item Ist die Annahme einer laminaren Strömungsform berechtigt?
\end{itemize}
 

\ifshowme
\antwort{
Umformen:
\begin{align*}
P &= \frac{8 \eta \cdot \Delta l}{\pi \cdot r^4} \cdot \dot{V}^2 \\
\dot{V} \cdot \Delta p &=  \frac{8 \eta \cdot \Delta l}{\pi \cdot r^4} \cdot \dot{V}^2 \\
 \Delta p &=  \frac{8 \eta \cdot \Delta l}{\pi \cdot r^4} \cdot \dot{V} \\
r^4 &=  \frac{8 \eta \cdot \Delta l}{\pi \cdot  \Delta p} \cdot \dot{V} \\
r &= \sqrt[4]{\frac{8 \eta \cdot \Delta l}{\pi \cdot  \Delta p} \cdot \dot{V}}
\end{align*}
Werte einsetzen:
\begin{align*}
r &= \sqrt[4]{\frac{8 \cdot 0,031 ~\text{kg/ms} \cdot 100~\text{m}}{\pi \cdot  0,1 \text{bar}} \cdot 13~\text{m}^3/\text{h}} \\
r &= \sqrt[4]{\frac{8 \cdot 0,031 ~\text{kg/ms} \cdot 100~\text{m}}{\pi \cdot  10.000 ~\text{kg/ms}^2} \cdot 0,00361~\text{m}^3/\text{s}}\\
r &\approx 0,041~\text{m} \\
d &\approx 0,082~\text{m} \\
d &\approx 8,2~\text{cm}
\end{align*}
Schauen ob Strömung laminar oder turbulent:
\begin{align*}
\text{Re}_{\text{krit}} = \frac{\rho d v}{\eta} = 2300
\end{align*}
$v$ ausrechnen:
\begin{align*}
v = \frac{\dot{V}}{A} \rightarrow v &= \frac{13~\text{m}^3/\text{h}}{\pi \cdot (0,041~\text{m})^2} \\
&\approx 2461,65~\text{m/h} \\
&\approx 0,684~\text{m/s} 
\end{align*}
Werte einsetzen:
\begin{align*}
\frac{850~\text{kg/m}^3 \cdot 0,082~\text{m} \cdot 0,684~\text{m/s}}{0,031~\text{kg/ms}} \approx 1537,9 \rightarrow {laminar}
\end{align*}


}
\pagebreak
\else
\fi


\subsection*{Aufgabe 2}
Wasser von 20 °C ($\eta$ = 10$^{-3}$ Pa s ; $\rho$ = 1.000 kg/m$^3$) fließt mit einem Volumenstrom von 600 m$^3$/h durch ein Stahlrohr ($\xi$ = 0,0113) mit der Nennweite 200 mm. \par 
\begin{itemize}
\item Wie groß ist der Druckverlust bei einer Leitungslänge von 1.000 m?

\ifshowme
\antwort{
$\xi$ von Stahlrohr $\approx$ 0,0113 \par 
$\overline{v}$ ausrechnen:
\begin{align*}
\overline{v} = \frac{\dot{V}}{A} \rightarrow \overline{v} &= \frac{600~ \m^3/\text{h}}{\pi \cdot (0,1 \mskip3mu \text{m})^2} \\
 &\approx 19.098,6~\text{m/h} \\
 &\approx 5,3~\text{m/s}  \\ \\
\end{align*}
Wir nutzen die Formel für Druckabfall im Rohr:
\begin{align*}
\Delta p &= \xi \cdot \frac{l}{d} \cdot \frac{1}{2} \rho \cdot \overline{v}^2\\
\end{align*}
Werte einsetzen:
\begin{align*}
\Delta p &= 0,0113 \cdot \frac{1.000~\text{m}}{0,2~\text{m}} \cdot \frac{1}{2} \cdot 1.000 ~\text{kg/m}^3 \cdot (5,3~\text{m/s})^2 \\ &\approx 793.542~\text{Pa}\\
&\approx 7,94 \text{bar}
\end{align*}

}
\else
\fi

\item Wenn die \splmt{Rohrreibungszahl} nicht angegeben wäre, wie könnten Sie diese herausfinden?

\ifshowme
\antwort{
Über zum Teil aufwendige Gleichungen (Colebrook, Nikuradse, etc.) oder aber über das Auslesen eines \textbf{Moody-Diagramms}.
}
\else
\fi

\item Welche Faktoren fließen in die \splmt{Rohrreibungszahl} ein?


\ifshowme
\antwort{
Die \textbf{Rohrreibungszahl} ist Abhängig von:
\begin{itemize}
\item[$\triangleright$] relativer Rauhigkeit $e$ (Quotient aus Wandrauhigkeit und Rohrinnendurchmesser $k/d$) 
\item[$\triangleright$] Reynoldszahl
\end{itemize}
}
\else
\fi
\end{itemize}
\marginnote[]{ \small 
Der dimensionslose \fplmt{Widerstandsbeiwert $\zeta$} beschreibt den Druckverlust speziell für einzelne Bauteile, während die dimensionslose \fplmt{\mbox{Rohrreibungszahl $\xi$}} (oder auch $\lambda$ in anderer Literatur) den kontinuierlichen Reibungsverlust entlang eines geraden Rohres beschreibt. \pp 
}[-4.75cm]





\subsection*{Aufgabe 3}
Der Volumenstrom von Bier, das durch ein Rohr fließt, beträgt 1,8 L/s. Der
Innendurchmesser des Rohres beträgt 3 cm. Das Bier hat eine Dichte von 1.100 kg/m$^3$.

\begin{itemize}
\item Berechnen Sie die durchschnittliche Geschwindigkeit von Bier in m/s und seinen Massendurchfluss in kg/s.
\item Wie hoch ist die Geschwindigkeit bei gleichem Volumenstrom, wenn ein anderes Rohr mit einem Durchmesser von 1,5 cm verwendet wird?
\end{itemize}

\ifshowme
\antwort{
$\dot{V} = 1,8 ~$L$/$s $ = 0,0018~$m$^3/$s \par 
$r=0,015~$m \pp
Massendurchfluss:
\begin{align*}
0,0018~\m^3/\s \cdot 1100~\kg/\m^3 = 1,98~\kg/\s
\end{align*}
Mittlere Geschwindigkeiten für versch. Durchmesser:
\begin{align*}
\overline{v} = \frac{\dot{V}}{A} \rightarrow \qquad \overline{v} &= \frac{0,0018~\m^3/s}{\pi \cdot (0,015~\text{m})^2} \qquad = 2,546~\m/\s \\
\overline{v} &= \frac{0,0018~\m^3/\s}{\pi \cdot (0,0075~\text{m})^2} \qquad = 10,19~\m/\s \\
\end{align*}
}
\pagebreak
\else
\fi


%\pagebreak


\subsection*{Aufgabe 4}
\tcbdocmarginnote[colback=unigelb!50,colframe=uniblau]{
Wiederholungsaufgabe
}
Ein Edelstahltank mit einem Durchmesser von 3 m enthält Wein. Der Tank wird bis zu
einer Höhe von 5 m befüllt. Eine Auslassöffnung mit einem Durchmesser von 10 cm wird
geöffnet, um den Wein abzulassen. \par Berechnen Sie die Austragsgeschwindigkeit des
Weins, unter der Annahme, dass der Fluss gleichmäßig und reibungslos ist. Wie lange
dauert es, bis der Tank vollständig entleert ist?


\ifshowme
\antwort{
Toricelli:
\begin{align*}
t = \sqrt{\frac{2}{g}} \cdot \frac{A_1}{A_2} \cdot (\sqrt{h_{\alpha}}-\sqrt{h_{\omega}})
\end{align*}
Werte einsetzen:
\begin{align*}
t = \sqrt{\frac{2}{9,81~\m/\s^2}} \cdot \frac{\pi \cdot (1,5~\m)^2}{\pi \cdot (0,05~\m)^2} \cdot \sqrt{5~\m} \qquad = 908,67~\s
\end{align*}
}
\pagebreak
\else
\fi



\subsection*{Aufgabe 5}
Zeichnen Sie die Fließ- und Viskositätskurve einer newton'schen Flüssigkeit, einer
strukturviskosen Flüssigkeit und einer dilatanten Flüssigkeit.
\ppbig


\marginnote[]{
\small Die \fplmt{Schubspannung $\tau$} gibt an, welche (horizontale) Kraft $F_R$ wir auf eine gegebene Fläche $A$ ausüben, um diese parallel zu verschieben. \par Die \fplmt{Scherrate $\dot{\gamma}$} gibt an, wie sich die Fließgeschwindigkeit in $x$-Richtung über die Höhe $z$ aufgrund der ausgeübten Kraft verändert.
}[-1.65cm]
\begin{marginfigure}[2.5cm]
\begin{tikzpicture}[scale=1.6]

\draw[miugray!50,anchor=south.west,fill=miugray!50](0,0) -- (0.5,0.75) -- (2,0.75) -- (1.5,0) -- cycle ;
\draw[dashed,thick](0,0)--(0.5,0.75);
\draw (0,0.3) to [bend left=45] node [above,xshift=1pt,scale=0.8] {$\gamma$} (0.2,0.3);
%rectangle (2.9cm,0.9cm);

\draw[->, miugray!50!black, dashed] (0.125,0.1875) -- (1.625,0.1875);
\draw[->, miugray!50!black, dashed] (0.25,0.375)-- (1.75,0.375);
\draw[->, miugray!50!black, dashed] (0.375,0.5625)-- (1.875,0.5625);

\draw[miugray!50,fill=miugray!50] (0.75, 0.23) rectangle (1.25,0.52);
\node[uniblau!75!black,scale=0.8](1)at(1,0.375){\bfseries \sffamily Fluid};
\draw[->](0,0)to(2.5,0)node[right]{$x$};
\draw[anchor=south.west,fill=miugray](0.5,0.9) rectangle node[above,yshift=5pt]{$A$} (2cm,0.75cm);
\draw[<->](0,0)to node[left]{$z$}(0,0.75);
\draw[->](2,0.85)to node[above]{$v$}(3,0.85);
\draw[->](2.5,0.7)to node[below]{$F_R$}(3.5,0.7);


\end{tikzpicture}
\caption{Modellvorstellung zur Bestimmung der dynamischen Viskosität; eine ebene Platte wird mit einer gewissen Kraft auf einem Flüssigkeitsfilm parallel zur Unterlage verschoben}
\end{marginfigure}


\begin{tikzpicture}[scale=3]
\node(a)at(0,0){
\begin{tikzpicture}[scale=2]
\node[align=center,scale=1](1)at(0.5,1.5){Newtonsch};
\draw (0,1)node[right]{$\tau$} -- (0,0) -- (1,0)node[right,scale=0.8]{$\dot{\gamma}$};
\draw[yshift=-1.5cm] (0,1)node[right]{$\eta$} -- (0,0) -- (1,0)node[right,scale=0.8]{$\dot{\gamma}$};
\ifshowme
{
\draw[black!80] (0,0) -- (0.9,0.9);
\draw[black!80,yshift=-1.5cm] (0,0.5) -- (0.9,0.5);
}
\else
\fi
\end{tikzpicture}
};
\node(b)at(1,0){
\begin{tikzpicture}[scale=2]
\node[align=center,scale=1](1)at(0.5,1.5){Scherverdünnend};
\draw (0,1)node[right]{$\tau$} -- (0,0) -- (1,0)node[right,scale=0.8]{$\dot{\gamma}$};
\draw[yshift=-1.5cm] (0,1)node[right]{$\eta$} -- (0,0) -- (1,0)node[right,scale=0.8]{$\dot{\gamma}$};
\ifshowme
{
\draw[black!80] (0,0) to [bend left=30] (0.9,0.9);
\draw[black!80,yshift=-1.5cm] (0,0.9) to [bend right=45] (0.9,0.3);
}
\else
\fi
\end{tikzpicture}
};
\node(c)at(2.1,0){
\begin{tikzpicture}[scale=2]
\node[align=center,scale=1](1)at(0.5,1.5){Scherverdickend};
\draw (0,1)node[right]{$\tau$} -- (0,0) -- (1,0)node[right,scale=0.8]{$\dot{\gamma}$};
\draw[yshift=-1.5cm] (0,1)node[right]{$\eta$} -- (0,0) -- (1,0)node[right,scale=0.8]{$\dot{\gamma}$};
\ifshowme
{
\draw[black!80] (0,0) to [bend right=30] (0.9,0.9);
\draw[black!80,yshift=-1.5cm] (0,0.3) to [bend right=30] (0.9,0.9);
}
\else
\fi
\end{tikzpicture}
};
\node[scale=0.9,rotate=90](fl)at(-0.75,0.45){Fließkurve};
\node[scale=0.9,rotate=90,align=left](vi)at(-0.75,-0.65){Viskositäts-\\kurve};
\end{tikzpicture}




\subsection*{Aufgabe 6}
\tcbdocmarginnote[colback=unigelb!50,colframe=uniblau]{
Wiederholungsaufgabe
}
In \figurename{~\ref{masse}} (auf der nächsten Seite) sehen Sie den fünfstufigen Herstellungsprozess eines Fertigprodukts. 
\begin{itemize}
\item Bestimmen Sie den Massenstrom von Produkt $P$
\item Bestimmen Sie den Massenstrom von Rücklauf $A$
\end{itemize}


\begin{figure}
%\resizebox{.5\textwidth}{!}{
\begin{tikzpicture}[font=\sffamily, >={Latex[length=6pt, width=5pt]}]

    % --- Definition der Stile für die Blöcke ---
    \tikzset{
        block/.style={rectangle, draw=black, minimum width=1.6cm, minimum height=1.4cm, align=center},
		blaublock/.style={block, fill=uniblau!50},
        grayblock/.style={block, fill=miugray!50},
        redblock/.style={block, fill=miugray}
    }

    % --- Positionierung der Blöcke ---
    \node[blaublock] (I)   at (0, 0)      {I};
    \node[redblock]  (II)  at (0, -2.5)   {II};
    \node[redblock]  (III) at (-1.5, -5.0){III};
    \node[redblock]  (IV)  at (1.5, -5.0) {IV};
    \node[grayblock] (V)   at (1.5, -7.5) {V};

    % --- Systemgrenze (gestrichelte Linie) ---
    \draw[dashed, semithick, uniblau] (-3.2, 1) rectangle (3.2, -8.7);
    \node[above left, inner sep=2pt] at (3.2, 1.2) {\color{uniblau}Systemgrenze};

    % --- Stoffströme (Pfeile & Beschriftungen) ---
    
    % Feed F
    \draw[->] (-4.8, 0) -- (-2.7, 0);
    \draw[->] (-2.7, 0) -- (I.west);
    \node[left, align=left, inner sep=2pt] at (-5, 0) {Feed F \\ 1000 kg/h\\15\% TS\\85\% Wasser};

    % Stream D
    \draw[->] (I.east) -- (4.2, 0);
    \node[right, align=left, inner sep=2pt] at (4.5, 0) {D \\ 150 kg/h\\0\% TS};

    % Stream B
    \draw[->] (I.south) -- (II.north) node[midway, right] {B};

    % Stream C, E und G (Aufteilung)
    \draw[->]  (II.south) -- (0, -3.75) node[midway, right] {C};
    \draw[-]  (0, -3.75) -- (-1.5, -3.75) node[midway, below,xshift=5pt] {E 10\% TS};
    \draw[->] (-1.5, -3.75) -- (III.north);
    \draw[-]  (0, -3.75) -- (1.5, -3.75) node[right] {G};
    \draw[->] (1.5, -3.75) -- (IV.north);

    % Stream A (Rückführung)
    \draw[-]  (III.west) -- (-2.7, -5.0);
    \draw[->] (-2.7, -5.0) -- (-2.7, 0);
    \node[right, align=left] at (-2.7, -2.5) {A\\5\% TS};

    % Stream R (Rückführung)
    \draw[-]  (IV.east) -- (2.7, -5.0) -- (2.7, -2.5);
    \draw[->] (2.7, -2.5) -- (II.east);
    \node[left] at (2.7, -3.75) {R};

    % Stream L
    \draw[->] (IV.south) -- (V.north) node[midway, right] {L};

    % Stream K
    \draw[->] (III.south) -- (-1.5, -9.0) node[below, align=left] {K\\450 kg/h\\20\% TS};

    % Stream Product P
    \draw[->] (V.south) -- (1.5, -9.0) node[below, align=left] {Produkt P\\80\% TS};

    % Stream W
    \draw[->] (V.east) -- (4.2, -7.5) node[right, align=left] {W\\Wasser};

\end{tikzpicture}
%}

\caption{Schematische Darstellung des Herstellungsprozesses. Abgeändert nach Singh und Heldman.}
\label{masse}
\end{figure}

\end{document}
